\chapter{Interconnects and arbitration}\label{ch:interconnects}

\whatyoufind{The three-part model every interconnect is built from; the bus family that the
shipped systems use, with its lanes, latencies and back-pressure; the three arbiters and the
one rule they share; the mesh and the single-switch NoC, with what they do and do not model
today; and where each of these is a plug-in.}

\section{Topology, controller, interfaces}

\figmed{Interconnect2.pdf}{An interconnect is a topology that creates a connection map and
instantiates one controller and the interfaces; the interfaces are the queues that hold
messages at each node; the controller moves messages between interfaces according to the map
and the configured latencies.}{fig:interconnect}{0.5}

An interconnect has two parts (\cref{fig:interconnect}). The \emph{topology}
(\code{InterconnectTopology}) is a set of message-holding interfaces plus a \emph{connection
map} describing who talks to whom. The \emph{controller} (\code{InterconnectController}) moves
messages between interfaces, and its arbiter resolves contention on shared links. Swapping the topology and controller pair yields a point-to-point
link, a unified, split or triple bus, a mesh or a NoC; swapping the arbiter yields FCFS,
round-robin or TDM. The controllers on either side see only interfaces, so nothing else in the
system changes.

Interconnects have their own clock: \key{m\_clk\_period} is 50 for every bus and mesh against
100 for the cores, so an interconnect steps twice per core cycle. The reason is resolution, not
speed: the periods are the units the \code{ClockManager} schedules in (\cref{ch:system-organization}),
and an interconnect that ticks between two core ticks can arbitrate a lane and deliver what it
carries within the same core cycle, so a request that wins the lane is at the LLC's queue one
core cycle after it was issued, and transfer times can be set in half core cycles. A transfer
occupies a link for \key{m\_request\_latency} (2) interconnect ticks for a request and
\key{m\_response\_latency} (5) for data: one core cycle and two and a half, the 2- and 5-tick
slots the visualizer draws. Both are configuration keys, so a slower link is a number, not a
clock change.

\section{The bus family}

\begin{figure}[htbp]\centering
\begin{adjustbox}{max width=\textwidth}% The three shipped interconnect shapes between the L1s and the LLC.
\begin{tikzpicture}[font=\sffamily\scriptsize, >=Latex,
  n/.style={draw=octrule, fill=octmem, rounded corners=2pt, minimum width=9mm, minimum height=5.5mm, font=\sffamily\scriptsize\bfseries},
  llc/.style={n, fill=octctl, draw=octamber, minimum width=12mm},
  sw/.style={draw=octpurple, fill=octio, circle, minimum size=7mm, font=\sffamily\scriptsize\bfseries},
  lane/.style={octpurple, thick},
  link/.style={octrule, thick},
  ttl/.style={font=\sffamily\small\bfseries, text=octink},
  sub/.style={font=\sffamily\scriptsize, text=octquiet, align=center, text width=44mm}]
  % ---- TripleBus ----
  \begin{scope}
    \node[ttl] at (18mm, 14mm) {TripleBus (shipped)};
    \foreach \i in {0,1,2,3} { \node[n] (b\i) at (\i*11mm, 6mm) {L1 \i}; \draw[link] (b\i.south) -- (\i*11mm, -2mm); }
    \foreach \y/\t in {-2/request, -5/response, -8/service} { \draw[lane] (-4mm, \y mm) -- (37mm, \y mm) node[right, text=octpurple] {\t}; }
    \node[llc] (bl) at (16.5mm, -16mm) {LLC}; \draw[link] (bl.north) -- (16.5mm, -8mm);
    \node[sub] at (18mm, -25mm) {three lanes, one arbiter each; every L1 sees every request (snooping); the service lane carries back-invalidations};
  \end{scope}
  % ---- Mesh ----
  \begin{scope}[xshift=62mm]
    \node[ttl] at (18mm, 14mm) {Mesh (\code{MultiCoreSystem\_Mesh})};
    \foreach \i/\x/\y in {0/0/6, 1/12/6, 2/24/6, 3/36/6} { \node[n] (m\i) at (\x mm, \y mm) {L1 \i}; }
    \node[llc] (ml) at (18mm, -10mm) {LLC};
    \foreach \a/\b in {m0/m1, m1/m2, m2/m3, m0/ml, m1/ml, m2/ml, m3/ml, m0/m2, m1/m3, m0/m3} { \draw[link] (\a) -- (\b); }
    \node[sub] at (18mm, -25mm) {\code{connection\_matrix} lists every node's neighbours; as shipped it is the complete graph, so every pair has a direct link with its own request and data channels};
  \end{scope}
  % ---- NoC ----
  \begin{scope}[xshift=124mm]
    \node[ttl] at (18mm, 14mm) {NoC (one switch)};
    \node[sw] (s) at (18mm, -3mm) {sw};
    \foreach \i/\x/\y in {0/0/6, 1/12/6, 2/24/6, 3/36/6} { \node[n] (n\i) at (\x mm, \y mm) {L1 \i}; \draw[link] (n\i) -- (s); }
    \node[llc] (nl) at (18mm, -14mm) {LLC}; \draw[link] (nl) -- (s);
    \foreach \i in {0,1,2,3} { \draw[link, octquiet, densely dotted] (n\i.south) to[bend right=12] (nl.north); }
    \node[sub] at (18mm, -25mm) {L1s link to the LLC directly (dotted) and to each other through switch 200, which forwards in the same cycle from a static next-hop table; at most one switch per message today};
  \end{scope}
\end{tikzpicture}
\end{adjustbox}
\caption{The three shapes that ship between the L1s and the LLC.}
\label{fig:topologies}
\end{figure}

\begin{description}
\item[Point-to-point] (\code{Point2PointController}) joins exactly two components, one message
  per direction per tick, no arbitration. \code{bus[1]} between the LLC and memory is one.
\item[Unified bus] (\code{UnifiedBusController}) carries requests and data on one shared link
  under one arbiter.
\item[Split bus] (\code{SplitBusController}) separates the request lane from the response lane,
  each with its own arbiter, so a long data transfer does not block the next request.
\item[TripleBus] (\code{TripleBus} with a \code{Split} controller; \cref{fig:topologies} left)
  adds a third, \emph{service} lane for the LLC's back-invalidations. It is what the shipped
  snooping systems use as \code{bus[0]}: every L1 sees every request on the request lane, data
  travels the response lane to its destinations only, and an \code{INV} on the service lane is
  broadcast to every interface including the LLC's own.
\end{description}

\paragraph{Back-pressure.} Each interface holds \key{buffers\_max\_size} messages (256 as
shipped). On the TripleBus \key{response\_reserve} (128) keeps that many entries of the response
buffers free for in-flight responses: the request lane stalls when the response buffer's free
space falls to the reserve, so responses --- which must always be able to drain --- can never
overflow. Delivery order at the receiving side is the subject of \cref{ch:cache-controller}.

\section{Arbiters}\label{sec:arbiters}

\begin{figure}[htbp]\centering
\begin{adjustbox}{max width=\textwidth}% The same three ready requests under FCFS, round-robin and TDM on one request lane
% (slot = 2 cycles). Core 0 has two requests ready at cycle 0, core 1 one.
\begin{tikzpicture}[font=\sffamily\scriptsize, >=Latex, x=9mm, y=1mm]
  \tikzset{slot/.style={draw=octrule, minimum height=5mm, anchor=west, inner sep=0, font=\sffamily\scriptsize\bfseries, text=white},
           c0/.style={slot, fill=octteal}, c1/.style={slot, fill=octamber}, idle/.style={slot, fill=octpaper, text=octquiet, draw=octquiet, densely dashed},
           rowl/.style={anchor=east, font=\sffamily\footnotesize\bfseries}}
  % axis
  \foreach \t in {0,...,10} { \draw[octquiet] (\t, 2) -- (\t, 0); \node[below, text=octquiet] at (\t, 0) {\t}; }
  \draw[octquiet] (0, 2) -- (10, 2);
  \node[text=octquiet] at (5, -6) {cycle (request-lane slot = 2 cycles)};
  % ready marks
  \node[anchor=east, text=octquiet] at (-0.3, 40) {ready at 0:};
  \node[c0, minimum width=9mm] at (0, 40) {C0 a};
  \node[c0, minimum width=9mm] at (1.1, 40) {C0 b};
  \node[c1, minimum width=9mm] at (2.2, 40) {C1};
  % FCFS
  \node[rowl] at (-0.3, 30) {FCFS};
  \node[c0, minimum width=18mm] at (0, 30) {C0 a};
  \node[c0, minimum width=18mm] at (2, 30) {C0 b};
  \node[c1, minimum width=18mm] at (4, 30) {C1};
  \node[anchor=west, text=octquiet] at (6.2, 30) {arrival order; C1 waits 4};
  % RR
  \node[rowl] at (-0.3, 20) {RR};
  \node[c0, minimum width=18mm] at (0, 20) {C0 a};
  \node[c1, minimum width=18mm] at (2, 20) {C1};
  \node[c0, minimum width=18mm] at (4, 20) {C0 b};
  \node[anchor=west, text=octquiet] at (6.2, 20) {one slot per core per round; C1 waits 2};
  % TDM
  \node[rowl] at (-0.3, 10) {TDM};
  \node[c0, minimum width=18mm] at (0, 10) {C0 a};
  \node[c1, minimum width=18mm] at (2, 10) {C1};
  \node[idle, minimum width=18mm] at (4, 10) {slot 2 idle};
  \node[idle, minimum width=18mm] at (6, 10) {slot 3 idle};
  \node[c0, minimum width=18mm] at (8, 10) {C0 b};
  \node[anchor=west, text=octquiet, text width=20mm, align=left] at (10.1, 10) {fixed slots; C0 b waits 8, but never more};
\end{tikzpicture}
\end{adjustbox}
\caption{Three requests ready at the same cycle on one request lane --- two from core 0, one
from core 1 --- under the three shipped arbiters.}
\label{fig:arbitration}
\end{figure}

An arbiter is a subclass of \code{Arbiter} implementing one method:

\begin{lstlisting}[style=cpp]
virtual bool elect(uint64_t cycle_number,
                   vector<vector<Message> *> &buffers,   // one TX buffer per interface
                   Message *out_msg) = 0;                // the winner, if any
\end{lstlisting}

It is called once per interconnect tick for each lane that is free, with every interface's TX
buffer, and names the message that gets the lane (\cref{fig:arbitration}):

\begin{description}
\item[FCFS] (\code{FCFSArbiter}) picks the oldest message across all buffers: minimal average
  wait, no bound on how long a core may be starved by a busier one.
\item[Round-robin] (\code{RRArbiter}) rotates a pointer over the cores and grants the pointed
  core's oldest message if it has one, otherwise the next core that does: every core gets at
  most one slot per round, so a core's wait is bounded by one slot per competitor.
\item[TDM] (\code{TDMArbiter}) gives each core a fixed slot in a fixed cycle, whether or not it
  has anything to send; an owner with nothing ready leaves its slot idle. The wait is bounded by
  the period regardless of what other cores do --- the property real-time analyses want --- at
  the cost of throughput. It is the shipped default on \code{bus[0]}
  (\file{SplitBusController.csv}).
\end{description}

An arbiter that grants by owner takes the owner's \emph{oldest} message across all buffers,
not the first buffer in interface order that holds one: a core's cache-to-cache supplies sit in a
different buffer from its own requests and would otherwise be passed over round after round.
The helper is \code{electOldestOwned} in \code{Arbiter}; a new arbiter should use it too.

The arbiter of a bus is a configuration key, \key{arbiter\_type}, in the controller's CSV or
overridden per instance: \code{-p "bus[0].interconnect\_controller.arbiter\_type(s)=RRArbiter"}.
The LLC's data port has its own arbiter under the same interface, \key{llc\_controller.arbiter\_type}.
Adding an arbiter is a subclass, a line in \file{src/Arbiters/CMakeLists.txt} and a name in the
controller's factory (\cref{ch:extending}).

\section{Mesh and NoC}

\code{MultiCoreSystem\_Mesh.csv} replaces the bus with a \code{Mesh}: one node per entry of
\key{component\_ids}, and \key{connection\_matrix[i]} listing node $i$'s neighbours, the first
entry being the node itself. As shipped (\file{configuration/Interconnect/Mesh.csv}) the matrix
is the complete graph on the four L1s and the LLC (\cref{fig:topologies} middle), so every pair
has a direct link:

\begin{lstlisting}[style=csv]
component_ids(vi),0,1,2,3,10
connection_matrix[0](vi),0,1,2,3,10    # node 0 (L1 0) links to L1 1, L1 2, L1 3 and the LLC
connection_matrix[1](vi),1,0,2,3,10
connection_matrix[2](vi),2,0,1,3,10
connection_matrix[3](vi),3,0,1,2,10
connection_matrix[4](vi),10,0,1,2,3    # the LLC links to every L1
\end{lstlisting}

A different shape is a different matrix, no code: a ring is each node listing its two
neighbours. Each link each link carries requests and data on separate channels with
their own latencies (2 and 5 ticks) and its own arbiter (FCFS by default), half-duplex, one
message in flight per latency window. It is the natural home of the directory protocol, which
sends point-to-point and never broadcasts; directory MSI runs the whole EEMBC suite on it.

The \code{NoC} is the mesh plus \emph{switch} nodes (\key{switch\_ids}, 200 as shipped). A
switch is an interface that forwards whatever it receives in the same cycle, with no latency or
arbitration of its own, from a static per-source next-hop table built at construction: every
direct neighbour maps to itself, and everything a neighbouring switch lists maps to that
switch. As shipped (\cref{fig:topologies} right) the L1s link to the LLC directly and to each
other through the switch. A message can cross at most one switch, because a switch's own table
knows only its direct neighbours.

\begin{notebox}[What the NoC does not model yet]
Three things are worth knowing before measuring on it. A destination the routing table does not
know is dropped silently and \code{pushMessage} returns true; the TX-side capacity check
compares the \emph{number of per-link queues}, a constant, against the limit, so there is no
back-pressure on that side and congestion exits instead of stalling; and a hop is logged as one
opaque bus event, so the report's bus columns and the visualizer cannot see the path. A
multi-hop network with pluggable routing, per-hop credits and virtual networks per message
class remains future work; the current single-switch NoC is the supported implementation.
\end{notebox}

\begin{wherebox}
\begin{itemize}[nosep,leftmargin=1.2em]
\item \file{header/Interconnect/}, \file{src/Interconnect/} --- \code{Bus}, \code{TripleBus} and their interfaces
\item \file{src/Interconnect/BusController.cpp}, \file{src/Interconnect/SplitBusController.cpp}, \file{src/Interconnect/TripleBusController.cpp}, \file{src/Interconnect/UnifiedBusController.cpp}, \file{src/Interconnect/Point2PointController.cpp} --- the bus controllers
\item \file{src/Interconnect/Mesh.cpp}, \file{src/Interconnect/MeshController.cpp}, \file{src/Interconnect/NoC.cpp}, \file{src/Interconnect/NoCController.cpp} --- the mesh and the NoC
\item \file{header/Arbiters/}, \file{src/Arbiters/} --- \code{Arbiter}, \code{FCFSArbiter}, \code{RRArbiter}, \code{TDMArbiter}
\item \file{configuration/Interconnect/} --- latencies, buffer depths, arbiter names, the mesh and NoC connection matrices (\cref{app:config-keys})
\end{itemize}
\end{wherebox}
